\section{Introduction}
\label{sec:intro}

Causal attribution --- determining which factors brought about a given outcome, and to
what degree --- is a central reasoning task in many important problems, such as
gauging the efficacy of medical treatments \citep{Shepherd2024}, pinpointing drivers
of economic growth \citep{CheshireMagrini2009}, root cause analysis in automated
systems \citep{Papageorgiou2022}, or assessing the impact of public policy
\citep{HeckmanPinto2022CausalPolicyAnalysis}. We may want to understand why an
adverse outcome occurred so as to prevent recurrence, or how a predicted positive
outcome could be brought about.

In the absence of scalable, formal causal attribution tools, analysts often resort to
ad-hoc methods: strong correlation, historical comparison, domain expertise, or
non-causal attribution like Granger causality \citep{Granger1969} and SHAP
\citep{lundberg2017unified}. Even researchers with an advanced command of
statistical and experimental methods may not possess a thorough understanding of how
to ascribe cause \citep{Shepherd2024}, and would benefit from scalable automation.
Here we present a tool for accurate and tractable causal attribution applicable to a
wide array of causal, probabilistic, machine-learned models.

What would a satisfactory causal attribution method look like? At minimum, it
should be (i) \textit{causally grounded}, tracking counterfactual dependence
rather than correlation; (ii) \textit{context-sensitive}, distinguishing
causes whose path was active in the factual world from those whose path was
blocked; (iii) \textit{broadly applicable}, handling stochastic models,
discrete and continuous variables, and finite (not just infinitesimal)
counterfactual changes; and (iv) \textit{tractable}, estimable from samples
without combinatorial enumeration over witness sets and alternatives. Existing
methods each forfeit one or more of these. Actual causality
\citep{actualCausalityHalpern} is grounded and context-sensitive, but neither
broadly applicable nor tractable in general, with no clear approximation
\citep{eiterComplexityResultsStructurebased2002}. Methods that \textit{do}
scale either give up causal grounding or are restricted in scope. SHAP
\citep{lundberg2017unified} grounds attribution in observational correlations,
so a feature can receive high attribution merely for being correlated with a
true cause; Causal SHAP \citep{heskes2020causal} addresses this to some extent,
but failure modes persist (Section~\ref{sec:shap_examples}).
\citet{butler2022differential}'s Differential Causal Effect (DCE) is causally
grounded but local: it relies on gradients, so requires differentiability and
does not capture effects from non-infinitesimal changes in potential causes.%
\footnote{DCE is, fairly, not designed for responsibility attribution:
\citeauthor{butler2022differential} target rates of change of treatment.}
Standard treatment-effect quantities, ATE and CATE, are causally grounded and
often tractable but lack context-sensitivity: every intervention propagates downstream
through the structural equations, overriding the mediator values that mark which
causal paths were active in the factual world
(Section~\ref{sec:motivations}).

We introduce \emph{probabilistic causal impacts} (\Ourapproach), an attribution
framework that combines the context-sensitivity of actual causality with the
scalability of probability-of-causation methods. \Ourapproach generalises Pearl's
probability of necessary and sufficient causation \citep{pearl2022probabilities} from
binary to arbitrary variable domains, and incorporates the \emph{witness} mechanism
from actual causality \citep{actualCausalityHalpern} via integration over a
distribution of witness sets rather than existential quantification.
Intuitively, a witness set records the variables held at their factual values while a
candidate cause is varied, so that causal paths inactive in the factual world stay
blocked rather than re-activating under intervention --- delivering the
context-sensitivity that ATE and CATE forgo. The effect is to
read actual causation --- otherwise an all-or-nothing verdict, and graded only
comparatively, by the relative normality of the witnessing contingency
\citep{Halpern2015-HALGCA} --- as a cardinal probabilistic magnitude on the same
structural information.
Several components have precedents --- expectations over
exogenous noise \citep{actualCausalityHalpern}, weighting by alternative-value
probabilities \citep{Beckers2015CombiningPC}, aggregation over witness sets
\citep{Beckers2015CombiningPC, actualCausalityHalpern} --- but no prior construction
combines such intuitively compelling components into a single tractable estimator
handling arbitrary variable domains.
Crucially, every quantity in the \Ourapproach definition is a standard interventional
or counterfactual expectation that causal probabilistic-programming systems such as
ChiRho%
\footnote{\url{https://basisresearch.github.io/chirho/explainable_sir.html}}
already evaluate automatically for a broad class of probabilistic causal models, so
\Ourapproach can be computed without bespoke per-model derivation.

\ourapproach presumes access to a trained probabilistic causal model from which
counterfactual outcomes can be sampled, a stronger requirement than methods that
work from observational data and a causal graph alone (e.g., the PN/PS/PNS bounds of
\citet{muellerpearl2022causes}, which return \emph{intervals} on the probability of
causation where \ourapproach, given a fully specified model, returns point values, or
Causal SHAP). Within that requirement, however, we place no strong restrictions on the form
of the model, distributions, or types of variables.
Such models can comprise standard machine learning components such as neural
network parametrized distributions or Gaussian processes.
Allowing neural components makes this far less restrictive than it first appears:
\citet{xia2022neural} show that neural causal models, being universal approximators,
are expressive enough to recover nonparametric counterfactual-identification results,
so a sufficiently flexible trained model forfeits no identification power available to
an explicitly specified structural model.
Nor do we require that counterfactual outcomes be point-identified: \ourapproach is an
expectation over the counterfactual distribution the model induces, so it is
well-defined whenever the model captures the relevant causal uncertainty, whether or
not that distribution is uniquely pinned down from data.

With this breadth of application, and a measure that surfaces its design choices ---
the variable selection distribution, alternative value distribution, and causal impact
function --- as explicit user-facing components rather than hidden defaults,
\ourapproach hopes to provide causal attribution and explanation tooling for a
wide array of users and applications.

\paragraph{Paper structure.}
The paper proceeds in four phases.
The first establishes motivation and formalism.
Section~\ref{sec:motivations} works through a running example
to surface what actual causality, the probability of necessity/sufficiency,
and average/conditional treatment effects each get right and miss, motivating
the design choices behind \ourapproach.
Section~\ref{sec:causal_impact} then defines the method: a kernel built from a
sufficiency check at the factual configuration and a necessity check against
alternatives, integrated over a distribution of suspect/witness sets, with
three components exposed as user choices: the variable selection distribution
$\Gamma$, the alternative-value distribution $\Delta$, and the causal impact
function.

The second phase positions \ourapproach against existing attribution methods.
Section~\ref{sec:shap_examples} compares \ourapproach with SHAP, Causal SHAP,
ATE, and CATE on two small examples (OBCB, with two binary features, and a
signal-with-mediation model), tabulating which desiderata each satisfies or fails,
including cases where SHAP and Causal SHAP attributions track correlation rather
than counterfactual dependence.
Section~\ref{sub:synthetic} reports a synthetic evaluation built around
overdetermination and undercutting, the patterns on which pre-AC accounts most
often diverge from intuition, and shows that \ourapproach recovers the literature's
verdicts.
Section~\ref{sec:DCE} contrasts \ourapproach with gradient-based attribution,
instantiated by Differential Causal Effect \citep{butler2022differential}, on a
continuous credit-limit example, showing where the two diverge and why
\ourapproach better tracks responsibility under unit rescaling and outside the
steep region of the response surface.

The third phase establishes the formal and empirical relationship between
\ourapproach and actual causality. Section~\ref{sec:relation_with_ac} shows that
under specific choices of $\Gamma$ and $\Delta$ the \ourapproach kernel is
positive on a configuration exactly when that configuration is an actual cause in
Halpern's sense, so aproximate AC verdicts can be read off the kernel without enumerating
witnesses.
Section~\ref{sub:pearl_actual_cause_prob} positions \ourapproach against Pearl's
probability of actual causation: on the canonical desert-traveller example,
\ourapproach recovers Pearl's within-scenario ranking with graded magnitudes,
and a weak-poison variant surfaces two additional responsibility intuitions
(path-reliability and cross-scenario reliability) that Pearl's binary indicator
structurally cannot express but \ourapproach's expectation does.
Section~\ref{sec:ac_benchmark} then benchmarks \ourapproach against explicit
actual-cause computation: exact subset enumeration times out at $17$ variables,
while the necessity-only \ourapproach estimator continues to roughly $73$--$145$
variables (still with relatively undemanding sample budget and  per-problem-size computation duration <=60 minutes) and sustains above a $0.85$
correct-attribution rate up to about $60$ variables, visiting several orders of
magnitude fewer cause--witness configurations.

The fourth phase pushes \ourapproach into settings where actual causality no
longer applies.
Section~\ref{sec:sir_benchmark} extends the benchmarking to a continuous-outcome
dynamical-SIR (Susceptible--Infected--Recovered) setting --- a regime where classical
actual causality returns only a binary, all-or-nothing verdict and cannot rank
competing causes. \ourapproach instead delivers a graded comparison, assigning the
lockdown policy roughly twice the responsibility of masking --- a distinction that
simple counterfactual (``but-for'') reasoning cannot draw.
Section~\ref{sec:avm} applies \ourapproach to a deployed automated valuation
model with machine-learned components trained on millions of points, a regime
where actual causality is intractable; SHAP returns qualitatively different
attributions, assigning virtually all responsibility to very few downstream variables
and effectively ignoring the rest.
Section~\ref{sec:discussion} closes with discussion.

\paragraph{Reproducibility.}
Every quantitative claim in this paper --- the worked-example numbers, the
benchmark tables, and the figures --- is reproduced in a set of companion
notebooks, rendered with their full outputs at
\url{https://rfl-urbaniak.github.io/explainable_paper/}. Throughout, we point to
the relevant notebook by name (for example \nb{obcb\_computations}); each name
links to its page on that site.



